\documentclass[10pt,a4paper]{amsart}
\usepackage[margin=19mm]{geometry}
\usepackage{amsmath,amssymb,mathtools}
\usepackage[scr=rsfs,cal=euler]{mathalfa}
\usepackage{xfrac}
\usepackage[unicode,hidelinks,bookmarks=false]{hyperref}
\newcommand{\ie}{i.e.\ }
\newcommand{\cf}{cf.\ }
\newcommand{\resp}{respectively\ }
\newcommand{\Subsection}{\subsection}
\providecommand{\nobreakdash}{\nobreak\hbox{-}}
\setlength{\emergencystretch}{2em}
\multlinegap0pt
\usepackage{tikz}
\usepackage[normalem]{ulem}
\usepackage{amssymb}
\usepackage{dsfont}
\usepackage{cancel}
\usepackage{braket}
\newtheorem{lem}{Lemma}
\newcommand		{\N}		{\mathbb N}			% Natural numbers
\newcommand		{\R}		{\RR}
\newcommand		{\RR}		{\mathbb R}			% real numbers
\newcommand		{\bangle}[1]	{\lt\langle #1\rt\rangle}
\newcommand		{\ceil}[1]		{\lt\lceil{#1}\rt\rceil}
\newcommand		{\indic}	{\mathds{1}}		% indicatrice
\newcommand		{\lt}			{\left}				%
\newcommand		{\rt}			{\right}			%
\newcommand		{\weight}[1]	{\bangle{#1}}	% <x>
\title{Commutator Estimates for Low-Temperature Fermi Gases}
\author{Jacky J. Chong, Laurent Lafleche, Jinyeop Lee, Chiara Saffirio}
\date{Source: arXiv:2604.02297v1 (CC BY 4.0)}
\begin{document}
\maketitle

\noindent\emph{Source and licence.} Commutator Estimates for Low-Temperature Fermi Gases. Version arXiv:2604.02297v1 (CC BY 4.0), distributed by arXiv under the Creative Commons Attribution 4.0 International licence, http://creativecommons.org/licenses/by/4.0/, as recorded in the arXiv licence metadata for this exact version and shown on the abstract page https://arxiv.org/abs/2604.02297v1. Full licence notice: \texttt{licenses/arxiv-2604.02297-CC-BY-4.0.txt}.\par\smallskip

\noindent\textit{Editorial scope.} Counted unit: the article's lem lem (source0.tex lines 1508--1518). The article's complete original proof of it is delivered with it (source0.tex lines 1520--1593, one \texttt{proof} environment of 74 lines). No mathematics is written, repaired or re-derived: the statement and the proof are the article's own lines, in the article's own notation, and no retained line is substituted or re-flowed.  No further source-local context is needed: every cross-reference of every retained line resolves inside the two delivered spans.  Nothing was omitted from any retained span, so the package carries no omission record.  No retained line contains a citation, so the wrapper carries no bibliography and no bibliography entry.  No retained line contains a figure environment, an image-inclusion command or drawing code, so no image is delivered and the package declares no asset dependency.  Editorial formatting only, in addition: an AMS \texttt{amsart} preamble that restates the article's own declarations and macros used by the retained lines, namely the environments \texttt{lem} on separate counters, together with the standard packages those lines need.  The retained environments print in this document as Lemma 1 in order of appearance, whereas the article prints the same items with the numbering of its own full document.  The complete native source of the article as distributed in the arXiv e-print for this exact version is bundled unchanged as \texttt{sources/197725/source0.tex}.
	\begin{lem}
		Let $\delta\in\R$, $\sigma\in\R_+^3$, and $p\geq 1$. Then there exists $C>0$, independent of $\delta, \sigma$, and $p$, such that 
		\begin{multline*}
			\Sigma:=\sum_{n\in\N_0^3:\, \sigma\cdot n \geq \delta} e^{-\sigma\cdot n +\delta}\, n_j^\frac{p}{2} \le \lt(\frac{p}{e}\rt)^{\frac{p}{2}}\frac{e^\delta}{\sigma_j^{p/2}} \prod^3_{i=1} \frac{\weight{\sigma_i}}{\sigma_i}\indic_{\{\delta< 0\}}
            \\
            + \frac{C}{\sigma_j^{p/2}}
			\lt(\lt( p^\frac{p}{2} + \lt(2\delta\rt)^\frac{p}{2} \rt)
			+ \lt(1+\frac{\delta}{\weight{\sigma_{j+1}}}\rt) \frac{\delta^{\frac{p}{2}+1}}{\weight{\sigma_j}} \rt) \prod_{k=1}^3 \frac{\weight{\sigma_k}}{\sigma_k}\indic_{\{\delta\ge 0\}}\,,
		\end{multline*}
		where $j+1$ can be replaced by $j+2$ and is modulo $3$.
	\end{lem}

	\begin{proof}
        Suppose $\delta<0$. Then we have 
        \begin{equation*}
            \Sigma=e^\delta\operatorname{Li}_{-\frac{p}{2}+1}(e^{-\sigma_j}) \prod_{\substack{i=1\\ i\ne j}}^3 (1-e^{-\sigma_i})^{-1}.
        \end{equation*}
        We have that 
        \begin{align*}
            \Sigma \le&\, \frac{2^{p/2}e^\delta}{\sigma_j^{p/2}} \lt(\sum^\infty_{n=1} \lt(\frac{\sigma_j n}{2}\rt)^{\frac{p}{2}} e^{-\sigma_j n}\rt)\prod_{\substack{i=1\\ i\ne j}}^3 (1-e^{-\sigma_i})^{-1} \\
            \le&\, \lt(\frac{p}{e}\rt)^{\frac{p}{2}}\frac{e^\delta}{\sigma_j^{p/2}} \lt(\sum^\infty_{n=1} e^{-\frac{\sigma_j n}{2}}\rt)\prod_{\substack{i=1\\ i\ne j}}^3 (1-e^{-\sigma_i})^{-1}\le \lt(\frac{p}{e}\rt)^{\frac{p}{2}}\frac{e^\delta}{\sigma_j^{p/2}} \prod^3_{i=1} \frac{\weight{\sigma_i}}{\sigma_i}\,,
        \end{align*}
        where the second inequality follows from the fact that for any $t,c\geq 0$, $t^c \leq \lt(c/e\rt)^c e^t$.

        
		Assume $\delta>0$ and take $j=1$. Then
		\begin{equation*}
			\Sigma 
            = e^{\delta}\sum_{n_1\geq 0} n_1^\frac{p}{2}\,e^{-\sigma_1 n_1} \sum_{n_2\geq 0} e^{-\sigma_2 n_2} \sum_{n_3 \geq \frac{(\delta-\sigma_1 n_1-\sigma_2 n_2)_+}{\sigma_3}} e^{-\sigma_3 n_3} ,
		\end{equation*}
        where $(x)_+:= \max(x, 0)$. Using the fact that for $r\in\R_+$ and $x\in(0,1)$, we have
		\begin{equation*}
			\sum_{k\geq r} x^k = \frac{x^{\ceil{r}}}{1-x} \leq \frac{x^r}{1-x} \, ,
		\end{equation*}
		then it follows that
		\begin{equation*}
			\Sigma \leq \frac{e^\delta}{1-e^{-\sigma_3}} \sum_{n_1\geq 0} n_1^\frac{p}{2}\,e^{-\sigma_1 n_1} \sum_{n_2\geq 0} e^{-\sigma_2 n_2 - (\delta-\sigma_1 n_1-\sigma_2 n_2)_+} .
		\end{equation*}
        
		Notice that 
        \begin{align*}
        \sigma_2 n_2 + (\delta - \sigma_1 n_1 - \sigma_2 n_2)_+=
            \begin{cases}
                \sigma_2 n_2 & \text{ if } \sigma_2 n_2 \geq \delta -\sigma_1 n_1
                \\
                \delta-\sigma_1 n_1 & \text{ else}
            \end{cases}.
        \end{align*}
        Hence, we have that
		\begin{align*}
			\sum_{n_2\geq 0} e^{-\sigma_2 n_2 - (\delta-\sigma_1 n_1-\sigma_2 n_2)_+} &= \sum_{n_2\geq \frac{(\delta-\sigma_1 n_1)_+}{\sigma_2}} e^{-\sigma_2 n_2} + \sum_{0\leq n_2 < \frac{(\delta-\sigma_1 n_1)_+}{\sigma_2}} e^{ - \delta+\sigma_1 n_1}
			\\
			&\leq \frac{e^{- (\delta-\sigma_1 n_1)_+}}{1-e^{-\sigma_2}} +  \frac{(\delta-\sigma_1 n_1)_+}{\sigma_2} \, e^{- \delta+\sigma_1 n_1} .
		\end{align*}
		So, it follows that
		\begin{equation*}
			\Sigma \leq \frac{J_1 + J_2}{\lt(1-e^{-\sigma_3}\rt)\lt(1-e^{-\sigma_2}\rt)}
		\end{equation*}
		where
		\begin{align*}
			J_1 &:= e^\delta \sum_{n_1\geq \delta/\sigma_1}  n_1^\frac{p}{2} e^{-\sigma_1 n_1}
			\\
			J_2 &:= \sum_{0\leq n_1 < \delta/\sigma_1}  n_1^\frac{p}{2} \lt(1 + \frac{\delta-\sigma_1 n_1}{\sigma_2}\lt(1-e^{-\sigma_2}\rt)\rt) .
		\end{align*}
		The second sum is bounded by
		\begin{equation*}
			J_2 \leq \lt(1+\frac{2\,\delta}{\weight{\sigma_2}}\rt) \lt(\frac{\delta}{\sigma_1}\rt)^{\frac{p}{2}+1} ,
		\end{equation*}
		which, again, follows from using the fact that $t^c \leq \lt(c/e\rt)^c e^t$. The first term is bounded by
		\begin{align*}
			J_1 &= \frac{2^{p/2}\,e^\delta}{\sigma_1^{p/2}}\sum_{n_1\geq \delta/\sigma_1} \lt(\tfrac{\sigma_1}{2}\,n_1\rt)^\frac{p}{2} e^{-\sigma_1 n_1}
			\\
			&\leq \frac{2^{p}\,e^\delta}{\sigma_1^{p/2}} \sum_{n_1\geq \delta/\sigma_1} \lt(\tfrac{\sigma_1 n_1 - \delta}{2}\rt)^\frac{p}{2} e^{-\sigma_1 n_1} + \lt(\tfrac{\delta}{2}\rt)^\frac{p}{2} e^{-\sigma_1 n_1}
			\\
			&\leq \frac{2^p e^\delta}{\sigma_1^{p/2}}\sum_{n_1\geq \delta/\sigma_1} \lt(\lt(\tfrac{p}{2e}\rt)^\frac{p}{2} e^{-\delta/2}\, e^{-\sigma_1 n_1/2} + \lt(\tfrac{\delta}{2}\rt)^\frac{p}{2} e^{-\sigma_1 n_1} \rt)
			\\
			&\leq \frac{e^\delta}{\sigma_1^{p/2}}\lt( p^\frac{p}{2}\,  \frac{e^{- \delta}}{1- e^{-\sigma_1/2}} + \lt(2\delta\rt)^\frac{p}{2} \frac{e^{- \delta}}{1- e^{-\sigma_1}} \rt) \leq \frac{2 p^\frac{p}{2} + \lt(2\delta\rt)^\frac{p}{2}}{\sigma_1^{p/2}\lt(1- e^{-\sigma_1}\rt)}\,,
		\end{align*}
		where we used the fact that $1 + e^{-\sigma_1/2} \leq 2$ to get the last inequality. Therefore, finally we have
		\begin{equation*}
			\Sigma \leq \frac{1}{\lt(1-e^{-\sigma_3}\rt)\lt(1-e^{-\sigma_2}\rt)}
			\lt(\frac{\lt( 2 p^\frac{p}{2} + \lt(2\delta\rt)^\frac{p}{2} \rt)}{\sigma_1^{p/2}\lt(1- e^{-\sigma_1}\rt)} 
			+  \lt(1+\frac{2\,\delta}{\weight{\sigma_2}}\rt) \lt(\frac{\delta}{\sigma_1}\rt)^{\frac{p}{2}+1} \rt)
		\end{equation*}
		and the result follows using the fact that $\frac{7}{8} \frac{t}{\weight{t}} \leq 1-e^{-t} \leq \frac{t}{\weight{t}}$.
	\end{proof}

\end{document}
